\section{Framework and Method}
\label{sec:framework}

We design a three-stage cross-domain trust management framework.
The first stage aligns heterogeneous trust scores into a shared semantic
space. The second stage detects domain-level trust-behavior drift by
comparing public semantic trust distributions with current and historical
behavioral observations. The third stage, which is the focus of this
section, attributes the detected domain-level drift to individual nodes.
The key design principle is to avoid directly thresholding QoS metrics.
Instead, we measure the residual between a node's trust-conditioned
behavioral expectation and its current behavior, and then quantify how much
the node contributes to the domain-level distributional drift.

\subsection{System Model and Notation}
\label{sec:notation}

Let $\mathcal{D}$ denote the set of trust domains. Each domain
$d \in \mathcal{D}$ contains a set of nodes $\mathcal{V}_d$.
After semantic alignment, each node $v \in \mathcal{V}_d$ has a unified
semantic trust score $s_v \in [0,1]$. The score $s_v$ represents the
public trust capability of node $v$ on a domain-independent semantic axis.

At epoch $t$, the system collects interaction observations and aggregates
them into node-level behavioral scores. We denote by $q_v(t) \in [0,1]$
the current observed behavior score of node $v$, and by $h_v(t) \in [0,1]$
its historical behavior score over a previous stable window. Let $n_v(t)$
be the number of observations supporting $q_v(t)$.

For each domain $d$, the second stage constructs three distributions:
\begin{itemize}
    \item $P_d(t)$: public semantic trust distribution induced by aligned
    trust records $\{s_v\}_{v \in \mathcal{V}_d}$;
    \item $Q_d(t)$: current behavioral distribution induced by
    $\{q_v(t)\}_{v \in \mathcal{V}_d}$;
    \item $H_d(t)$: historical behavioral distribution induced by
    $\{h_v(t)\}_{v \in \mathcal{V}_d}$.
\end{itemize}
The domain-level drift score is
\begin{equation}
    \Delta_d(t) =
    \frac{
    \mathrm{JS}(P_d(t), Q_d(t)) +
    \mathrm{JS}(P_d(t), H_d(t)) +
    \mathrm{JS}(Q_d(t), H_d(t))
    }{3},
    \label{eq:domain-drift}
\end{equation}
where $\mathrm{JS}(\cdot,\cdot)$ is the Jensen-Shannon divergence.
If $\Delta_d(t)$ exceeds a domain-level threshold, the third stage is
activated to localize the nodes responsible for the drift.

\subsection{Node-Level Trust Residual Attribution}
\label{sec:node-attribution}

\paragraph{Trust-conditioned behavioral expectation.}
For each node $v$, we first estimate the behavior that should be expected
from it at epoch $t$. The expectation combines three sources of evidence:
the aligned semantic trust score, the node's own historical behavior, and
the historical behavior of nodes in the same semantic bucket.

Let $\mathcal{B}(v)$ denote the set of historical observations from nodes
in the same domain and the same semantic score bucket as $v$. We define
the bucket baseline as
\begin{equation}
    b_v(t) = \frac{1}{|\mathcal{B}(v)|}
    \sum_{u \in \mathcal{B}(v)} h_u(t).
    \label{eq:bucket-baseline}
\end{equation}
When the bucket has insufficient history, we fall back to the domain-level
historical mean, and then to $s_v$ if no historical observation is available.

The trust-conditioned expectation of node $v$ is
\begin{equation}
    e_v(t) =
    \alpha s_v + \beta h_v(t) + \gamma b_v(t),
    \quad \alpha + \beta + \gamma = 1.
    \label{eq:expectation}
\end{equation}
In our implementation, we use $\alpha=0.3$, $\beta=0.5$, and $\gamma=0.2$
by default. This weighting gives priority to the node's own historical
behavior while still preserving the semantic trust commitment.

\paragraph{Positive residual.}
The node-level residual is the gap between expected behavior and current
behavior:
\begin{equation}
    r_v(t) = e_v(t) - q_v(t).
    \label{eq:raw-residual}
\end{equation}
Only positive residuals are treated as malicious evidence, because behavior
better than expectation does not indicate trust degradation. We further
weight the residual by observation support:
\begin{equation}
    R_v(t) =
    \min\left(1, \frac{n_v(t)}{n_{\min}}\right)
    \cdot
    \max(0, r_v(t)).
    \label{eq:weighted-residual}
\end{equation}
This prevents nodes with too few observations from dominating the
attribution result.

\paragraph{Leave-one-out drift contribution.}
A high residual indicates that a node behaves worse than expected, but it
does not necessarily explain the domain-level drift. Therefore, we also
measure how much node $v$ contributes to $\Delta_d(t)$.
Let $\Delta_d^{(-v)}(t)$ be the domain drift recomputed after removing
node $v$ from the current behavioral distribution $Q_d(t)$. We define
the contribution of node $v$ as
\begin{equation}
    C_v(t) =
    \max\left(0, \Delta_d(t) - \Delta_d^{(-v)}(t)\right).
    \label{eq:contribution}
\end{equation}
If removing $v$ reduces the drift, then $v$ is a contributor to the
domain-level anomaly. If the drift remains unchanged, the node is not a
major explanation for the domain alarm.

\paragraph{Attribution score and residual state.}
We combine the local residual and the leave-one-out contribution into an
attribution score:
\begin{equation}
    A_v(t) =
    \lambda R_v(t) +
    (1-\lambda)\cdot
    \min\left(1, \frac{C_v(t)}{\eta}\right),
    \label{eq:attribution-score}
\end{equation}
where $\lambda$ controls the balance between local residual and
distributional contribution, and $\eta$ normalizes the contribution term.
We use $\lambda=0.7$ by default.

To avoid reacting to isolated fluctuations, we maintain a one-sided CUSUM
state:
\begin{equation}
    S_v(t) =
    \max\left(0, S_v(t-1) + A_v(t) - \kappa\right),
    \label{eq:cusum}
\end{equation}
where $\kappa$ is a slack term for benign variation. Node $v$ is marked
as suspicious if $S_v(t) > \tau_{\mathrm{node}}$.

\paragraph{Attack interpretation.}
The residual trajectory provides a simple interpretation of common attacks.
For a normal node, $S_v(t)$ remains close to zero. For a grayhole node,
$R_v(t)$ is repeatedly positive and $S_v(t)$ gradually crosses the
threshold. For a blackhole node, $q_v(t)$ is close to zero and
$S_v(t)$ increases quickly. We classify a suspicious node as blackhole
when $q_v(t) \leq \theta_{\mathrm{black}}$; otherwise it is classified
as grayhole. Discrimination attacks can be handled by applying the same
residual computation separately for each source domain and comparing the
resulting group states, although the minimum core design only requires
the aggregate residual attribution above.

\begin{algorithm}[t]
\caption{Node-Level Residual Attribution}
\label{alg:node-attribution}
\begin{algorithmic}[1]
\Require Domain $d$, semantic scores $\{s_v\}$, current observations
$\{q_v(t), n_v(t)\}$, historical observations $\{h_v(t)\}$,
domain drift $\Delta_d(t)$
\Ensure Ranked node attribution results
\ForAll{$v \in \mathcal{V}_d$}
    \State Compute semantic-bucket baseline $b_v(t)$
    \State $e_v(t) \gets \alpha s_v + \beta h_v(t) + \gamma b_v(t)$
    \State $R_v(t) \gets \min(1,n_v(t)/n_{\min})\cdot
    \max(0,e_v(t)-q_v(t))$
    \State Recompute $\Delta_d^{(-v)}(t)$ by removing $v$ from $Q_d(t)$
    \State $C_v(t) \gets \max(0,\Delta_d(t)-\Delta_d^{(-v)}(t))$
    \State $A_v(t) \gets \lambda R_v(t) +
    (1-\lambda)\min(1,C_v(t)/\eta)$
    \State $S_v(t) \gets \max(0,S_v(t-1)+A_v(t)-\kappa)$
    \If{$S_v(t) > \tau_{\mathrm{node}}$}
        \If{$q_v(t) \leq \theta_{\mathrm{black}}$}
            \State $\mathrm{label}(v) \gets \textsc{Blackhole}$
        \Else
            \State $\mathrm{label}(v) \gets \textsc{Grayhole}$
        \EndIf
    \Else
        \State $\mathrm{label}(v) \gets \textsc{Normal}$
    \EndIf
\EndFor
\State \Return nodes ranked by $S_v(t)$ and $A_v(t)$
\end{algorithmic}
\end{algorithm}

\subsection{Optional Randomized Confirmation}
\label{sec:randomized-confirmation}

The residual attribution stage is sufficient for the minimum design.
However, when the attribution score falls into an ambiguous region, the
system can run a lightweight randomized confirmation procedure. The goal
is not to create artificial QoS stress, but to separate node-specific
misbehavior from shared environmental fluctuation.

The verifier selects a suspicious node $v$ and a small set of control
nodes $\mathcal{U}_v$ from the same domain and semantic bucket. It then
uses a secret key $K$ to generate hidden bits
\begin{equation}
    z_i = \mathrm{PRF}(K, v, i) \bmod 2.
\end{equation}
If $z_i=0$, a normal-looking request is sent to $v$; if $z_i=1$, the same
request distribution is used for a randomly selected control node
$u \in \mathcal{U}_v$. The target nodes only observe ordinary traffic.
The hidden bits are used only by the verifier to assign requests to
treatment and control groups. A differential residual CUSUM can then be
used:
\begin{equation}
    Z_v(i) =
    \max\left(0, Z_v(i-1) +
    \left(y_u(i)-y_v(i)\right) - \kappa_a\right),
\end{equation}
where $y_v(i)$ and $y_u(i)$ are the observed compliance scores for the
suspicious and control nodes. If $Z_v(i)>\tau_a$, the node-specific
misbehavior is confirmed. This step is optional and is only invoked for
ambiguous cases; the main detection pipeline remains passive and
lightweight.

\subsection{Discussion}
\label{sec:discussion}

The proposed node-level attribution is different from direct QoS anomaly
detection. A node is not flagged merely because its current QoS is low.
Instead, the system asks whether the current behavior is lower than what
is expected from the node's aligned semantic trust, own history, and peer
bucket baseline, and whether the node explains the domain-level
distributional drift. This makes the third stage a natural continuation
of the first two stages: semantic alignment defines comparable trust
capability, domain-level distributional detection raises an alarm, and
node-level residual attribution localizes the entities responsible for the
alarm.
